Formalization is in progress. The companion uses Mathlib's determinant-one real matrices and its actual upper-half-plane action. Its declarations cover the fractional formula, denominator nonvanishing, the imaginary-part identity on both nonreal half-planes, the difference formula on $\mathbb H$, the center $\{\pm I\}$, the action kernel, and equality of that kernel with the center. Lean's scalar-action notation represents the manuscript's $g\cdot z$ (also abbreviated $gz$); the type abbreviation for the real group is $\operatorname{SL}_2(\mathbb R)$. The Section 1.2 companion proves the individual scaling and translation actions and upper-half-plane transitivity. Its shared proof applies the scaling and translation formulas successively, matching the edited orbit argument; the formal product expresses composition of those actions.

The difference proof puts the fractions over a common denominator, factors the numerator as $(ad-bc)(z-w)$, and uses determinant one, following the edited prose. The imaginary-part proof specializes Mathlib's division formula, whose numerator is obtained by multiplication by the conjugate denominator; $\operatorname{normSq}$ is the manuscript's squared complex modulus. The center proof reuses Mathlib's general scalar-center theorem and specializes the scalar square to one. The prose instead computes commutation with $a(2)=\operatorname{diag}(\sqrt2,1/\sqrt2)$ and $n(1)$. The kernel proof evaluates the fixed-point equation at $i$ and $2i$, while the prose uses the polynomial's infinitely many zeros. These different routes prove the same center and kernel statements.

The shared derivative theorem, exposed in the Section 1.3 companion, now proves the complex derivative of the fractional formula.

The global action on the Riemann sphere, its determinant-one rescaling over $\mathbb C$, its three-orbit classification (including the lower half-plane and the boundary), the quotient-group action and its faithfulness, and the positive-determinant extension still require their formal counterparts. The imaginary-part formula is not by itself a formalization of the orbit theorem. Mathlib's action of the full real general linear group uses conjugation for negative determinants; this companion restricts to determinant one, where its action is precisely the holomorphic fractional action of the prose.

The cocycle and central-character compatibility are deferred to Chapter 3 or later modular-forms material. Neither is a current prose claim in this section. The proved shared-library cocycle lemma is retained for reuse. The shared function called \texttt{factor} denotes $cz+d$; its use in these formulas does not introduce an additional mathematical assumption or a cocycle claim in this section.

The recorded machine checks for the implemented portions are in \texttt{lean/audits/chapter01-release-check-2026-11-12/record.json}. Their source hashes identify the current checked snapshot. Neither successful compilation nor the axiom audit completes the missing targets or Steven's correspondence review.
